\BourbakiExerciseGroup

\begin{enumerate}
\def\labelenumi{\arabic{enumi})}
\item
  Let A be a principal ideal domain, let M be a torsion A-module, let x, y be two elements of M, and let Aα and Aβ be the annihilators of x, y respectively; if δ is a gcd of α and β, show that the annihilator Aγ of \(x + y\) is such that γ divides the lcm αβ/δ of α and β, and is a multiple of αβ/δ². When α and β are distinct powers \(\pi^{\mu}\) , \(\pi^{\nu}\) of the same irreducible element π of A, show that \(\gamma = \pi^{\sup(\mu, \nu)}\) .
\item
  Let A be a principal ideal domain, let \(\pi\) be an irreducible element of A and let M be a \(\pi\) -primary module; show that for all \(x \in M\) and all \(\alpha \in A\) not a multiple of \(\pi\) , there exists a unique \(y \in M\) such that \(x = \alpha y\) (use the Bezout identity); put \(y = \alpha^{-1}x\) . Deduce that, if \(A_{(\pi)}\) denotes the principal ideal domain consisting of elements \(\beta/\alpha\) of the field of fractions of A, such that \(\alpha, \beta \in A\) and \(\alpha\) is not a multiple of \(\pi\) (VII, p.~48, Ex. 4), then there exists a unique \(A_{(\pi)}\) -module structure on M such that restriction of the ring of operators to A gives the original A-module structure on M. The \(A_{(\pi)}\) -module thus defined is said to be canonically associated to M; its submodules are identical to the A-submodules of M.
\item
  Let A be a principal ideal domain.
\end{enumerate}

\begin{enumerate}
\def\labelenumi{\alph{enumi})}
\item
  An A-module M is injective (II, p.~389, Ex. 11) if and only if, for all \(x \in \mathbf{M}\) and all \(\alpha \neq 0\) in A, there exists \(y \in \mathbf{M}\) such that \(x = \alpha y\) ; the A-module M is then said to be divisible. If E is an arbitrary A-module, then the sum of all the divisible submodules of E is the greatest divisible submodule \(\Delta(E)\) of E. We have \(\Delta(F) \subset \Delta(E)\) for every submodule F of E.
\item
  Let M be a divisible A-module, let \(\pi\) be an irreducible element of A and let \(x_{0}\) be an element of M whose annihilator has the form \(A\pi^{n}\) (n \textgreater{} 0 an integer). Define a sequence \((x_{s})\) of elements of M recursively by the conditions \(x_{s} = \pi x_{s+1}\) for each integer \(s \geqslant 0\) . Show that the submodule of M generated by the sequence \((x_{s})\) is isomorphic to the \(\pi\) -primary component \(U_{\pi}\) of the module U = K/A, where K is the field of fractions of A. c) Show that \(U_{\pi}\) is a divisible module, and that every submodule of \(U_{\pi}\) other than \(U_{\pi}\) and 0 is cyclic and generated by the class mod A of some power of \(\pi\) ; deduce that \(U_{\pi}\) is indecomposable (II, p.~393, Ex. 21).
\end{enumerate}

\(d\) ) Show that every indecomposable divisible \(\mathbf{A}\) -module is isomorphic either to \(\mathbf{K}\) or to one of the modules \(\mathrm{U}_{\pi}\) (use \(a\) ) and \(b)\) ).

\begin{enumerate}
\def\labelenumi{\alph{enumi})}
\setcounter{enumi}{4}
\tightlist
\item
  Show that every divisible A-module is isomorphic to the direct sum of a family of indecomposable divisible submodules (apply Zorn's Lemma to families of indecomposable divisible submodules whose sums are direct).
\end{enumerate}

\begin{enumerate}
\def\labelenumi{\arabic{enumi})}
\setcounter{enumi}{3}
\tightlist
\item
  A family \((x_{i})\) of nonzero elements of a module \(\mathbf{M}\) over a principal ideal domain \(\mathbf{A}\) is said to be pseudofree if every relation of the form \(\sum_{i} \alpha_{i} x_{i} = \beta y\) ( \(y \in \mathbf{M}, \alpha_{i} \in \mathbf{A}\) ,
\end{enumerate}

\(\beta \in \mathbf{A}\) ) implies the existence of a family of elements \(\alpha_{i}^{\prime} \in \mathbf{A}\) such that \(\alpha_{i}x_{i} = \beta \alpha_{i}^{\prime} x_{i}\) for each index \(i\) . The relation \(i \neq k\) then implies \(x_{i} \neq x_{k}\) ; the set of elements of a pseudofree family is called a pseudofree subset of \(\mathbf{M}\) .

\begin{enumerate}
\def\labelenumi{\alph{enumi})}
\item
  Let \(\pi\) be an irreducible element of A. Show that, in an A-module with annihilator \(\mathbf{A}\pi^n\) ( \(n \geqslant 1\) ), any element with annihilator \(\mathbf{A}\pi^n\) is pseudofree (use Bezout's identity).
\item
  Let \((x_{i})\) be a pseudofree family in an A-module M, and let N be the submodule of M generated by this family. Show that for every element \(\dot{x}\) of M/N there exists an element \(x \in \dot{x}\) whose annihilator is equal to the annihilator of \(\dot{x}\) in M/N. If also \(\dot{x}\) is pseudofree in M/N, show that the family consisting of the \((x_{i})\) and x is pseudofree in M.
\item
  Deduce from a) and b) that if the A-module M has nonzero annihilator, then it is a direct sum of cyclic modules (reduce to the case of a \(\pi\) -primary module, and apply Zorn's Lemma to pseudofree families of elements of M) (cf.~VII, p.~19, Th. 2).
\end{enumerate}

\begin{enumerate}
\def\labelenumi{\arabic{enumi})}
\setcounter{enumi}{4}
\item
  Let M be a finitely generated torsion module over a principal ideal domain A. Show that M is a module of finite length; consequently every nonempty set of submodules of M has a maximal element and a minimal element.
\item
  Let M be a torsion module over a principal ideal domain A, which is not finitely generated, but all of whose proper submodules are finitely generated. Show that there exists an irreducible element \(\pi\) of A such that M is isomorphic to the \(\pi\) -primary component \(U_{\pi}\) of the module U = K/A (Ex. 3, c)). (First of all show that M cannot have more than one nonzero \(\pi\) -primary component, and hence is a \(\pi\) -primary module for some irreducible element \(\pi \in A\) . Then show that M = \(\pi M\) by observing that M/ \(\pi M\) could not be finitely generated if it were nonzero; but M/ \(\pi M\) is a vector space over the field A/ \(\pi A\) , and would contain a proper subspace which was not finitely generated.)
\item
  Let M be a module over a principal ideal domain A; a submodule N of M is said to be pure if, for all \(\alpha \in A\) , we have \(N \cap (\alpha M) = \alpha N\) .
\end{enumerate}

\begin{enumerate}
\def\labelenumi{\alph{enumi})}
\item
  The submodule N of M is pure if and only if, for all \(\alpha \in A\) and every element \(\dot{x} \in M/N\) whose annihilator is \(A\alpha\) , there exists an element \(x \in \dot{x}\) whose annihilator is \(A\alpha\) .
\item
  Show that the torsion submodule of M is pure. A family \((x_{t})\) of elements of M is pseudofree (Ex. 4) if and only if the sum of the submodules \(Ax_{t}\) is direct and is a pure submodule of M. Show that if P and Q are submodules of M such that \(P \cap Q\) and \(P + Q\) are pure then P and Q are pure. If Q is a pure submodule of M and \(P \supset Q\) a submodule of M, then P is pure in M if and only if P/Q is pure in M/Q.
\item
  Show that every submodule N of M which is a direct factor of M is pure. The same is true for the union of a filtered family of direct factors of M. Conversely, if N is a pure submodule of M, and M/N is a direct sum of cyclic submodules, show that N is a direct factor of M (use a)). Give examples of direct factors P and Q of M such that \(P \cap Q\) (resp. \(P + Q\) ) is not a pure submodule (for the first example, take \(M = (\mathbf{Z}/4\mathbf{Z}) \oplus (\mathbf{Z}/2\mathbf{Z})\) , and for the second, take \(M = Z^{2}\) ).
\item
  Let N be a pure submodule of M, whose annihilator Aα is nonzero. If f is the natural homomorphism from M onto M/αM, show that f restricts to an isomorphism from N onto f(N), and that f(N) is a pure submodule of M/αM (use Bezout's identity). Deduce (with the help of c) and Ex. 4, c)) that \(f(N)\) is a direct factor in M/ \(\alpha M\) , and that N is a direct factor in M. In particular, if the torsion submodule of a module M has nonzero annihilator (which is the case when it is finitely generated), then it is a direct factor in M (cf.~VII, p.~20, Cor. 1 to Th. 2).
\item
  Let \(p\) be a prime number; let \(\mathbf{N}\) denote the direct sum of the cyclic \(\mathbf{Z}\) -modules \(\mathbf{Z} / (p^n)\) ( \(n \geqslant 1\) ), and let \(e_n\) denote the element of \(\mathbf{N}\) whose coordinates are all zero except for the \(n\) -th, which is equal to the class of 1 mod \(p^n\) . Let \(\mathbf{M}\) be the submodule of the \(\mathbf{Z}\) -module \(\mathbf{N} \times \mathbf{Q}\) generated by the elements \((e_n, p^{-n})\) . Show that the torsion submodule \(\mathbf{P}\) of \(\mathbf{M}\) is not a direct factor in \(\mathbf{M}\) (notice that every element of \(\mathbf{M} / \mathbf{P}\) is divisible by \(p^n\) for all \(n\) , and that no element of \(\mathbf{M}\) has this property except for multiples of \((0, 1)\) ).
\end{enumerate}

\begin{enumerate}
\def\labelenumi{\arabic{enumi})}
\setcounter{enumi}{7}
\tightlist
\item
  Let A be a principal ideal domain, let \(\pi\) be an irreducible element of A and let M be a \(\pi\) -primary module. Then an element \(x\) of M is said to have height \(n\) if \(x \in \pi^n\mathbf{M}\) and \(x \notin \pi^{n+1}\mathbf{M}\) . If \(x \in \pi^n\mathbf{M}\) for all \(n \geqslant 1\) then \(x\) is said to have infinite height.
\end{enumerate}

\begin{enumerate}
\def\labelenumi{\alph{enumi})}
\item
  Show that M is divisible if and only if all its elements have infinite height (Ex. 3).
\item
  Let x be an element of M of height n, such that \(\pi x = 0\) ; show that if \(x = \pi^{n}y\) then the cyclic submodule Ay of M is a direct factor of M (notice that Ay is a pure submodule, and apply Ex. 7, d)).
\item
  Show that a nondivisible \(\pi\) -primary module M always contains at least one cyclic nonzero direct factor (if \(x \in \mathbf{M}\) has finite height, and there exist integers \(m\) such that \(\pi^m x \neq 0\) has infinite height, then let \(s\) be the least such integer; write \(\pi^s x = \pi y\) , where \(y\) has height strictly greater than \(\pi^{s-1}x\) , and apply \(b\) ) to \(y - \pi^{s-1}x\) ). Deduce that any indecomposable \(\pi\) -primary module is either cyclic or isomorphic to a module \(\mathrm{U}_{\pi}\) (Ex. 3). d) Show that if every pure submodule of the \(\pi\) -primary module M is a direct factor in M, then M is a direct sum of a divisible module P and a \(\pi\) -primary module Q with nonzero annihilator. (Reduce to the case where M contains no nonzero divisible submodule, and argue by contradiction. Use b) to show that there would exist an infinite sequence \((x_k)\) of elements of height 1 in M such that: 1) the annihilator of \(x_k\) is \(\mathrm{A}\pi^{n_k}\) , where the sequence \((n_k)\) of integers is strictly increasing; and 2) the sum of the \(\mathrm{A}x_k\) is direct, and the sum of any finite number of these submodules is a direct factor in M. Then show that the submodule N of M generated by the elements \(x_k - \pi^{n_{k+1} - n_k}x_{k+1}\) is pure, but not a direct factor of M, arguing as in Ex. 7, e)).
\end{enumerate}

\begin{enumerate}
\def\labelenumi{\arabic{enumi})}
\setcounter{enumi}{8}
\tightlist
\item
  \begin{enumerate}
  \def\labelenumii{\alph{enumii})}
  \tightlist
  \item
    Let M be a \(\pi\) -primary module, and let \(M_{0}\) be the submodule of M consisting of elements of infinite height in M. Show that the \(\pi\) -primary module \(M/M_{0}\) contains no nonzero element of infinite height.
  \end{enumerate}
\end{enumerate}

\begin{enumerate}
\def\labelenumi{\alph{enumi})}
\setcounter{enumi}{1}
\tightlist
\item
  Let \(p\) be a prime number, and let \(E\) be the direct sum of the \(p\) -groups \(E_n = \mathbf{Z} / (p^n)\) ( \(n \geqslant 1\) ); let \(e_n\) denote the residue class of \(1 \mod p^n\) in \(E_n\) . Let \(H\) be the submodule of \(E\) generated by the elements \(e_1 - p^{n-1}e_n\) for all integers \(n \geqslant 2\) . Let \(M\) be the quotient module \(E / H\) . Show that the submodule \(M_0\) of elements of infinite height in \(M\) is the submodule \((E_1 + H) / H\) , which is isomorphic to \(E_1\) , and show that there is no nonzero element in \(M_0\) which has infinite height in \(M_0\) .
\end{enumerate}

¶ 10) a) Let M be a π-primary module with no nonzero elements of infinite height and let N be a pure submodule of M generated by a maximal pseudofree family of elements of M (Ex. 4 and 7, b)). Show that M/N is a divisible module (argue by contradiction, applying Ex. 8, c), Ex. 7, c), Ex. 7, b), and finally Ex. 4, c)).

\begin{enumerate}
\def\labelenumi{\alph{enumi})}
\setcounter{enumi}{1}
\item
  Consider the metrisable topology \(\mathcal{T}\) on M defined by taking as a basis of neighbourhoods of 0 the submodules \(\pi^n\mathbf{M}\) for all integers \(n\geqslant 0\) (Gen.~Top., III, p.~5). Then the quotient M/P of M by a submodule P is divisible if and only if P is everywhere dense in M under the topology \(\mathcal{T}\) .
\item
  Let \(P_{1}\) , \(P_{2}\) be two submodules of M which are pure, everywhere dense in M under the topology T, and direct sums of cyclic submodules. Show that \(P_{1}\) and \(P_{2}\) are isomorphic (notice that \(P_{1}/\pi^{n}P_{1}\) is isomorphic to \(M/\pi^{n}M\) for all n).
\item
  Let \(\hat{\mathbf{M}}\) be the completion of \(\mathbf{M}\) with respect to the topology \(\mathcal{T}\) (Gen.~Top., III, p.~24). Show that \(x \mapsto \alpha x\) is a strict morphism from \(\hat{\mathbf{M}}\) onto \(\alpha \hat{\mathbf{M}}\) for all \(\alpha \in \mathbf{A}\) , that \(\pi^n \hat{\mathbf{M}}\) is an open and closed submodule, and the closure of \(\pi^n \mathbf{M}\) , for all \(n\) , and that the quotient modules \(\hat{\mathbf{M}} / \pi^n \hat{\mathbf{M}}\) , \(\mathbf{M} / \pi^n \mathbf{M}\) are isomorphic.
\item
  Let \(\bar{\mathbf{M}}\) be the torsion submodule of \(\hat{\mathbf{M}}\) ; show that \(\bar{\mathbf{M}}\) is a \(\pi\) -primary module, that \(\bar{\bar{\mathbf{M}}} = \bar{\mathbf{M}}\) , and that \(\mathbf{M}\) is a pure submodule of \(\bar{\mathbf{M}}\) .
\item
  Suppose the \(\pi\) -primary module M is the direct sum of a family \((E_{\alpha})\) of cyclic submodules. Then T is discrete if and only if the annihilator of M is 0. In this case, let E be the product module \(\prod_{\alpha} E_{\alpha}\) and let F be the submodule of E consisting of elements all but
\end{enumerate}

countably many of whose coordinates are zero; show that \(\hat{\mathbf{M}}\) can then be identified (as a non-topological module) with the submodule of \(F\) consisting of \((x_{\alpha})\) such that, for each integer \(n\) , there are only finitely many indices \(\alpha\) for which \(x_{\alpha} \notin \pi^{n} E_{\alpha}\) . Then show that the four modules \(\mathbf{M}\) , \(\hat{\mathbf{M}}\) , \(\bar{\mathbf{M}}\) and \(F\) are distinct. Deduce that \(\bar{\mathbf{M}}\) is not a direct sum of cyclic modules (use \(e\) ), and that it is not a direct sum of indecomposable modules (notice that \(\bar{\mathbf{M}}\) contains no elements of infinite height, and apply Ex. 8, c)).

\begin{enumerate}
\def\labelenumi{\arabic{enumi})}
\setcounter{enumi}{10}
\tightlist
\item
  Let M be a \(\pi\) -primary module and let N be a submodule of M whose nonzero elements all have height \(\leqslant h\) in M.
\end{enumerate}

\begin{enumerate}
\def\labelenumi{\alph{enumi})}
\item
  Let P be a pure submodule of M such that \(P(\pi) \subset N(\pi)\) and \(P(\pi) \neq N(\pi)\) . Let a be an element of \(N(\pi)\) not belonging to \(P(\pi)\) , and let \(x_{0}\) be an element of \(a + P(\pi)\) whose height k in M is as large as possible; let \(x_{0} = \pi^{k} y_{0}\) . Show that the sum \(P + Ay_{0}\) is direct, and is a pure submodule of M.
\item
  Let B be a pseudofree subset of M such that, if P is the pure submodule generated by B, then \(\mathrm{P}(\pi) \subset \mathrm{N}(\pi)\) . Show that there exists a pseudofree subset \(C \supset B\) of M such that, if Q is the pure submodule of M generated by C, then \(\mathrm{Q}(\pi) = \mathrm{N}(\pi)\) (apply Zorn's Lemma, using a)).
\end{enumerate}

\begin{enumerate}
\def\labelenumi{\arabic{enumi})}
\setcounter{enumi}{11}
\tightlist
\item
  A \(\pi\) -primary module M is a direct sum of cyclic submodules if and only if there exists a pseudofree subset H in M such that, if N is the pure submodule of M generated by H, then \(M(\pi) = N(\pi)\) .
\end{enumerate}

\begin{enumerate}
\def\labelenumi{\alph{enumi})}
\setcounter{enumi}{1}
\item
  The module M is a direct sum of cyclic submodules if and only if M is the union of an increasing sequence \((P_{k})\) of submodules, such that the heights in M of the nonzero elements in each \(P_{k}\) are bounded. (To see that the condition is sufficient, use a) and Ex. 11, b).
\item
  Deduce from \(b)\) that if \(M\) is a direct sum of cyclic submodules, then every submodule of \(M\) is a direct sum of cyclic submodules.
\item
  Let M be a \(\pi\) -primary module with no elements of infinite height, having a countable system of generators ( \(x_{k}\) ). Show that M is a direct sum of cyclic submodules (if
\end{enumerate}

\(P_{k}\) is the submodule generated by \(x_{i}\) for \(i \leqslant k\) , then apply b) to the increasing sequence \((\mathrm{P}_{k})\) , using Ex. 5 applied to the decreasing sequence \((\mathrm{P}_{k} \cap \pi^{j}\mathrm{M})\) (for fixed k)).

\begin{enumerate}
\def\labelenumi{\alph{enumi})}
\setcounter{enumi}{4}
\tightlist
\item
  Let M be a \(\pi\) -primary module having a countable system of generators. Then M is a direct sum of indecomposable submodules if and only if the submodule N of elements of infinite height in M is pure (use d) and Ex. 7, c)).
\end{enumerate}

\begin{enumerate}
\def\labelenumi{\arabic{enumi})}
\setcounter{enumi}{12}
\tightlist
\item
  Let A be a principal ideal domain, let \(\pi\) be an irreducible element of A, and let M be a \(\pi\) -primary module over A. Let I be the set of ordinals \(\alpha\) (Set Theory, III, p.~225, Ex. 14) such that the cardinality of the set of ordinals \(< \alpha\) is at most equal to that of M. Define submodules \(\mathbf{M}^{(\alpha)}\) of M for all \(\alpha \in I\) by transfinite induction as follows: take \(\mathbf{M}^{(0)} = \mathbf{M}\) ; take \(\mathbf{M}^{(\alpha + 1)} = \pi \mathbf{M}^{(\alpha)}\) ; and if \(\alpha\) is a limit ordinal take \(\mathbf{M}^{(\alpha)}\) to be the intersection of all the \(\mathbf{M}^{(\beta)}\) for \(\beta < \alpha\) .
\end{enumerate}

\begin{enumerate}
\def\labelenumi{\alph{enumi})}
\item
  Show that there exists a least ordinal \(\tau \in I\) such that \(M^{(\tau + 1)} = M^{(\tau)}\) (the terminal ordinal of M). Show that \(M^{(\tau)}\) admits a complement in M, and is the direct sum of submodules isomorphic to \(U_{\pi}\) (Ex. 3). Say that M is reduced if \(M^{(\tau)} = 0\) .
\item
  Let M be a reduced \(\pi\) -module. For all \(x \in \mathbf{M}\) , show that there exists a greatest ordinal \(\gamma(x) \leqslant \tau\) such that \(x \in \mathbf{M}^{(\gamma(x))}\) (in particular, \(\gamma(0) = \tau\) ). Generalising the definition in Ex. 8, we say that \(\gamma(x)\) is the height of \(x\) in M. Show that if \(\gamma(x) < \gamma(y)\) , then \(\gamma(x + y) = \gamma(x)\) , and that if \(\gamma(x) = \gamma(y)\) , then \(\gamma(x + y) \geqslant \gamma(x)\) . We say that \(x\) is proper with respect to a submodule N of M if for all \(y \in \mathbf{N}\) the height \(\gamma(x + y)\) is at most equal to \(\gamma(x)\) ; show that then \(\gamma(x + y) = \inf(\gamma(x), \gamma(y))\) , for all \(y \in \mathbf{N}\) .
\item
  Let \(\alpha < \tau\) be an ordinal, and let \(\tilde{\mathbf{M}}^{(\alpha)}\) be the submodule of \(\mathbf{M}^{(\alpha)}\) consisting of those \(x\) such that \(\pi x = \pi y\) for some \(y \in \mathbf{M}^{(\alpha + 1)}\) . For those \(y \in \mathbf{M}^{(\alpha + 1)}\) having this property, the elements of the form \(u = x - y\) form a class mod \(\mathbf{M}(\pi) \cap \mathbf{M}^{(\alpha + 1)}\) in \(\mathbf{M}(\pi) \cap \mathbf{M}^{(\alpha)}\) ; if \(\varphi\) is the canonical homomorphism from \(\mathbf{M}(\pi) \cap \mathbf{M}^{(\alpha)}\) onto the quotient of this module by \(\mathbf{M}(\pi) \cap \mathbf{M}^{(\alpha + 1)}\) , then show that the map \(\psi\) which sends \(x\) to \(\varphi(u)\) is a homomorphism with kernel \(\mathbf{M}^{(\alpha + 1)}\) , and consequently that \(\tilde{\mathbf{M}}^{(\alpha)} / \mathbf{M}^{(\alpha + 1)}\) is isomorphic to \((\mathbf{M}(\pi) \cap \mathbf{M}^{(\alpha)}) / (\mathbf{M}(\pi) \cap \mathbf{M}^{(\alpha + 1)})\) .
\item
  Let N be a submodule of M, and let \(x \in \tilde{\mathbf{M}}^{(\alpha)}\) be an element of height \(\alpha\) ; show that x is proper with respect to N if and only if the corresponding elements u = x - y of \(\mathbf{M}(\pi) \cap \mathbf{M}^{(\alpha)}\) (cf.~c)) are proper with respect to N. There exists such an element x if and only if the image of \(\mathbf{N} \cap \tilde{\mathbf{M}}^{(\alpha)}\) under the homomorphism \(\psi\) is distinct from \((\mathbf{M}(\pi) \cap \mathbf{M}^{(\alpha)}) / (\mathbf{M}(\pi) \cap \mathbf{M}^{(\alpha + 1)})\) .
\end{enumerate}

\begin{enumerate}
\def\labelenumi{\arabic{enumi})}
\setcounter{enumi}{13}
\tightlist
\item
  Let M be a reduced \(\pi\) -primary module. In the notation of Ex. 13, for each ordinal \(\alpha < \tau\) we can consider \((\mathbf{M}(\pi) \cap \mathbf{M}^{(\alpha)}) / (\mathbf{M}(\pi) \cap \mathbf{M}^{(\alpha + 1)})\) as a vector space over the field \(F_{\pi} = A / A\pi\) ; let \(c(\alpha)\) denote the dimension of this vector space, and call \(c(\alpha)\) the multiplicity of the ordinal \(\alpha\) in the module M (cf.~VII, p.~24).
\end{enumerate}

\begin{enumerate}
\def\labelenumi{\alph{enumi})}
\item
  Show that if M and N are two \(\pi\) -primary modules each of which is a direct sum of cyclic submodules (which implies that \(\tau \leqslant \omega\) , the first infinite ordinal), then M and N are isomorphic if and only if they have the same terminal ordinal \(\tau\) and the multiplicity of every integer \(n < \tau\) is the same in M and in N.
\item
  Give an example of two \(\pi\) -primary modules with no elements of infinite height, which are not isomorphic, but in which each integer n has the same multiplicity (cf.~Ex. 10, f). c) Let M and N be two, reduced \(\pi\) -primary modules, each having a countably infinite system of generators. Show that if M and N have the same terminal ordinal \(\tau\) , and if each ordinal \(\alpha < \tau\) has the same multiplicity in M and in N, then M and N are isomorphic (« Ulm-Zippin Theorem »). (Given two submodules \(S \subset M\) and \(T \subset N\) , an isomorphism \(h\) from \(S\) onto \(T\) is said to be indicial if, for all \(x \in S\) , the height of \(x\) in \(M\) is equal to the height of \(h(x)\) in \(N\) . Define an indicial isomorphism from \(M\) onto \(N\) by induction, and by alternating the rôles played by \(M\) and \(N\) , in such a way as to reduce the problem to the following: given two finitely generated submodules \(S \subset M\) and \(T \subset N\) , an indicial isomorphism \(h\) from \(S\) onto \(T\) , and an element \(x \in M \cap \complement S\) , extend \(h\) to an indicial isomorphism from \(S + Ax\) onto a submodule of \(N\) containing \(T\) . Show that this can be reduced to the case where \(\pi x \notin S\) ; using Ex. 5, choose an element \(w_1\) in \((S + Ax) \cap \complement S\) with greatest possible height \(\alpha\) in \(M\) , then choose an element \(w_2\) in \(S \cap M^{(\alpha)}\) such that the height of \(\pi(w_1 + w_2)\) in \(M\) is greatest possible; show that \(w = w_1 + w_2\) is proper with respect to \(S\) (Ex. 13, b)). Let \(z = h(\pi w) \in T\) ; show that, if there exists an element \(t\) of \(N\) , of height \(\alpha\) in \(N\) , proper with respect to \(T\) , and such that \(\pi t = z\) , then an extension \(\bar{h}\) of \(h\) of the desired form can be obtained by putting \(\bar{h}(w) = t\) . It remains to prove the existence of the element \(t\) . If
\end{enumerate}

\[
\gamma (z) = \gamma (\pi w) = \alpha + 1 <   \tau ,
\]

show that there exist elements \(u\) in \(\mathbf{N}^{(\alpha)} \cap \complement \mathbf{T}\) such that \(z = \pi u\) , and that \(t\) may be taken to be any such element. If on the other hand \(\gamma(z) = \gamma(\pi w) > \alpha + 1\) or else if \(\gamma(z) = \gamma(\pi w) = \alpha + 1 = \tau\) , use Ex. 13, \(d\) ) to prove the existence of the element \(t\) , noticing that (in the notation of Ex. 13) the image under \(\psi\) of \(S \cap \tilde{\mathbf{M}}^{(\alpha)}\) in \((\mathbf{M}(\pi) \cap \mathbf{M}^{(\alpha)}) / (\mathbf{M}(\pi) \cap \mathbf{M}^{(\alpha + 1)})\) is finite dimensional over the field \(F_{\pi}\) )

* d) Obtain from c) a new proof of Ex. 12, d) (use also Th. 1 of VII, p.~14). *

\begin{enumerate}
\def\labelenumi{\arabic{enumi})}
\setcounter{enumi}{14}
\tightlist
\item
  \begin{enumerate}
  \def\labelenumii{\alph{enumii})}
  \tightlist
  \item
    Let A be a principal ideal domain, let M be a divisible A-module (Ex. 3), and let N be a torsion A-module. Show that \(\mathbf{M} \otimes_{\mathbf{A}} \mathbf{N} = 0\) .
  \end{enumerate}
\end{enumerate}

\begin{enumerate}
\def\labelenumi{\alph{enumi})}
\setcounter{enumi}{1}
\tightlist
\item
  Let M and N be two \(\pi\) -primary modules with no elements of infinite height; show that \(M \otimes_{A} N\) is a direct sum of cyclic submodules (use a) and Ex. 10, a)).
\end{enumerate}

